% ===================================================================================
% BAB 4: PANDUAN PENGGUNAAN e-SPPO UNTUK PEMILIH -- SURAT SUARA ELEKTRONIK (SSE)
% Berisi panduan lengkap bagi pemilih untuk menggunakan aplikasi SSE.
% File ini akan di-include oleh main.tex
% ===================================================================================

\part{Panduan Penggunaan e-SPPO untuk Pemilih -- Surat Suara Elektronik (SSE)}
\label{part:penggunaan_sse}
Bab ini ditujukan khusus bagi Anda, para pemilih. Di sini, Anda akan menemukan panduan lengkap tentang cara menggunakan aplikasi Surat Suara Elektronik (SSE) untuk memberikan hak suara Anda secara digital. Prosesnya dirancang agar sederhana, cepat, dan aman.

% -----------------------------------------------------------------------------------
% Section: Pengenalan Surat Suara Elektronik (SSE)
% -----------------------------------------------------------------------------------
\section{Pengenalan Surat Suara Elektronik (SSE)}
\label{sec:pengenalan_sse}
Surat Suara Elektronik (SSE) adalah antarmuka berbasis web yang berfungsi sebagai "kertas suara digital" Anda dalam sistem e-SPPO. Melalui SSE, Anda dapat melihat daftar kandidat, mempelajari visi dan misi mereka, serta memberikan pilihan Anda dengan beberapa klik saja.

\subsection{Keunggulan SSE bagi Pemilih}
\begin{itemize}
	\item \textbf{Mudah Diakses dan Digunakan:} Dengan antarmuka yang modern dan responsif, SSE dapat diakses dari berbagai perangkat (komputer, laptop, tablet, ponsel) dan sangat mudah digunakan.
	\item \textbf{Proses Cepat:} Alur penggunaan yang sederhana memungkinkan Anda memberikan suara hanya dalam hitungan menit, mengurangi waktu antrean di TPS.
	\item \textbf{Kerahasiaan Terjamin:} Pilihan Anda direkam secara anonim. Identitas Anda dienkripsi dengan teknologi \textbf{AES-256-CBC}, sehingga tidak dapat dihubungkan langsung dengan suara yang Anda berikan.
	\item \textbf{Akurasi Terjamin:} Sistem menghitung suara secara otomatis, menghilangkan risiko kesalahan hitung yang biasa terjadi pada metode manual.
\end{itemize}

% -----------------------------------------------------------------------------------
% Section: Cara Kerja dan Keamanan SSE
% -----------------------------------------------------------------------------------
\section{Cara Kerja dan Keamanan SSE}
\label{sec:cara_kerja_sse}
Memahami cara kerja SSE dapat memberikan Anda keyakinan bahwa suara Anda aman dan tercatat dengan benar.

\subsection{Alur Proses Teknis}
Saat Anda menggunakan SSE, berikut adalah proses teknis yang terjadi di belakang layar, sesuai dengan kode programnya:
\begin{enumerate}
	\item \textbf{Validasi Kredensial (\texttt{login.php}):} Ketika Anda memasukkan Nama Akun dan Kata Sandi, sistem akan mencocokkan \textit{hash} kata sandi yang Anda masukkan dengan yang tersimpan di basis data menggunakan fungsi \texttt{password\_verify()}. Sistem juga akan memeriksa apakah status pemilihan sedang berlangsung dan apakah akun Anda belum digunakan untuk memilih.
	\item \textbf{Pencatatan Kehadiran:} Jika login berhasil, sistem akan langsung menyimpan data kehadiran Anda (ID unik pemilih, jenis kelamin, dan konstituensi) ke dalam tabel \texttt{data\_kehadiran}. Proses ini terjadi \textbf{sebelum} Anda memberikan suara, sehingga memisahkan data kehadiran dari data pilihan.
	\item \textbf{Penyimpanan Sesi:} Informasi penting seperti ID unik, nama, jenis kelamin, dan \textit{Initialization Vector} (IV) untuk enkripsi disimpan sementara dalam sesi (\texttt{\$\_SESSION}) untuk digunakan pada proses selanjutnya.
	\item \textbf{Proses Perekaman Suara (\texttt{ballot\_process.php}):} Ini adalah bagian paling krusial.
	\begin{itemize}
		\item \textbf{Transaksi Atomik:} Seluruh proses penyimpanan suara dijalankan dalam sebuah transaksi basis data. Ini berarti semua langkah harus berhasil; jika satu langkah gagal, seluruh proses akan dibatalkan (\textit{rollback}) untuk mencegah data yang tidak konsisten.
		\item \textbf{Penguncian Data:} Sistem mengunci baris data Anda di tabel pemilih (\texttt{FOR UPDATE}) untuk mencegah upaya pemilihan ganda secara bersamaan (\textit{race condition}).
		\item \textbf{Enkripsi Identitas:} ID unik Anda dienkripsi menggunakan kunci rahasia dan IV unik milik akun Anda dengan algoritma \textbf{AES-256-CBC}. Hasil enkripsi inilah yang disimpan bersama suara Anda, bukan ID asli.
		\item \textbf{Pembuatan ID Suara Unik:} Sistem menghasilkan sebuah ID unik untuk surat suara Anda yang digabungkan dari ID terenkripsi Anda dan ID kandidat yang dipilih.
		\item \textbf{Penyimpanan Suara:} Data suara (ID suara, nomor urut kandidat, ID pemilih terenkripsi, dll.) disimpan ke tabel \texttt{data\_suara}.
		\item \textbf{Pembaruan Status:} Status Anda di tabel \texttt{data\_pemilih} diubah menjadi `SUDAH\_MEMILIH`.
		\item \textbf{Commit Transaksi:} Jika semua langkah sukses, perubahan disimpan secara permanen.
	\end{itemize}
	\item \textbf{Penghancuran Sesi:} Setelah suara berhasil direkam, sesi Anda akan langsung dihancurkan (\texttt{session\_destroy()}) untuk keamanan. Anda akan otomatis keluar dari sistem.
\end{enumerate}

% -----------------------------------------------------------------------------------
% Section: Ketentuan Penggunaan SSE
% -----------------------------------------------------------------------------------
\section{Ketentuan Penggunaan SSE}
\label{sec:ketentuan_sse}
\begin{itemize}
	\item \textbf{Satu Akun, Satu Suara:} Setiap akun pemilih hanya dapat digunakan untuk memberikan suara satu kali.
	\item \textbf{Jaga Kerahasiaan Akun:} Jangan pernah membagikan Nama Akun dan Kata Sandi Anda kepada siapa pun.
	\item \textbf{Gunakan Perangkat yang Aman:} Jika memilih dari TPS, gunakan hanya Perangkat Akses Surat Suara Elektronik (PASSE) yang disediakan. Jangan mengutak-atik perangkat tersebut.
	\item \textbf{Laporkan Kendala:} Jika terjadi galat atau masalah teknis, segera laporkan kepada panitia pemilihan yang bertugas.
\end{itemize}

% -----------------------------------------------------------------------------------
% Section: Panduan Langkah Demi Langkah
% -----------------------------------------------------------------------------------
\section{Panduan Langkah Demi Langkah}
\label{sec:langkah_penggunaan_sse}
Ikuti langkah-langkah berikut untuk memberikan suara Anda melalui SSE.

\subsection{Langkah 1: Mengakses Halaman Login}
Buka peramban web pada perangkat Anda, lalu masukkan alamat URL SSE yang diberikan oleh panitia. Anda akan disambut oleh halaman login yang menampilkan nama kegiatan pemilihan yang sedang berlangsung.
\begin{figure}[h]
	\centering
	\includegraphics[width=1\textwidth]{SSE_1.png}
	\caption{Halaman Login Aplikasi SSE.}
	\label{fig:sse_login}
\end{figure}

\subsection{Langkah 2: Masuk (Login) ke Akun Anda}
Masukkan \textbf{Nama Akun Pemilih} dan \textbf{Kata Sandi} Anda pada kolom yang tersedia, lalu klik tombol \textbf{Masuk}.
\begin{figure}[h]
	\centering
	\includegraphics[width=1\textwidth]{SSE_2.png}
	\caption{Mengisi Kredensial Akun Pemilih.}
	\label{fig:sse_login_filled}
\end{figure}

\subsection{Langkah 3: Memilih Kandidat}
Setelah berhasil login, Anda akan masuk ke halaman surat suara utama.
\begin{enumerate}
	\item \textbf{Pelajari Kandidat:} Luangkan waktu untuk melihat foto, nama, serta visi dan misi setiap kandidat yang ditampilkan.
	\item \textbf{Tentukan Pilihan:} Klik pada tombol radio (lingkaran) atau pada area kartu kandidat pilihan Anda. Tanda centang biru akan muncul untuk menandakan pilihan Anda.
\end{enumerate}
\begin{figure}[h]
	\centering
	\includegraphics[width=1\textwidth]{SSE_4.png}
	\caption{Memilih Salah Satu Kandidat.}
	\label{fig:sse_selecting_candidate}
\end{figure}

\subsection{Langkah 4: Mengirim dan Mengonfirmasi Pilihan}
\begin{enumerate}
	\item Setelah yakin, klik tombol biru \textbf{Kirim Suara Saya}.
	\item Sebuah kotak dialog konfirmasi akan muncul. Ini adalah kesempatan terakhir Anda untuk memeriksa kembali pilihan.
	\item Klik \textbf{Ya - Kirim Suara} untuk melanjutkan, atau \textbf{Tidak - Kembali ke Surat Suara} untuk batal.
\end{enumerate}
\begin{figure}[h]
	\centering
	\includegraphics[width=1\textwidth]{SSE_5.png}
	\caption{Kotak Dialog Konfirmasi Pilihan.}
	\label{fig:sse_confirm_vote}
\end{figure}
\importantbox{Setelah Anda menekan tombol "Ya - Kirim Suara", pilihan Anda bersifat \textbf{final} dan tidak dapat diubah kembali.}

\subsection{Langkah 5: Selesai dan Keluar (Logout)}
Jika suara berhasil dikirim, Anda akan melihat halaman konfirmasi akhir.
\begin{figure}[h]
	\centering
	\includegraphics[width=1\textwidth]{SSE_6.png}
	\caption{Halaman Konfirmasi Akhir.}
	\label{fig:sse_vote_recorded}
\end{figure}
Klik tombol \textbf{Selesai \& Kembali ke Login} untuk keluar sepenuhnya dari akun Anda.
\warningbox{Langkah ini \textbf{wajib} dilakukan, terutama jika Anda menggunakan perangkat umum di TPS, untuk melindungi privasi dan mencegah penyalahgunaan akun.}

% -----------------------------------------------------------------------------------
% Section: Rekomendasi dan Informasi Tambahan
% -----------------------------------------------------------------------------------
\section{Rekomendasi untuk Pengguna SSE}
\label{sec:rekomendasi_sse}
\begin{itemize}
	\item \textbf{Gunakan Koneksi Stabil:} Pastikan Anda terhubung ke jaringan internet atau LAN yang stabil saat memilih untuk menghindari kegagalan pengiriman suara.
	\item \textbf{Jangan Terburu-buru:} Baca semua informasi dengan teliti sebelum membuat pilihan dan mengirimkan suara.
	\item \textbf{Jika Ragu, Bertanya:} Jangan ragu untuk bertanya kepada panitia jika ada hal yang tidak Anda mengerti selama proses pemilihan.
\end{itemize}

\section{Informasi Lainnya}
\label{sec:info_lainnya_sse}
\begin{itemize}
	\item \textbf{Pesan Galat (\textit{Error Messages}):} Jika Anda menemui pesan galat seperti "Nama Akun atau Kata Sandi salah" atau "Akun sudah digunakan", ikuti petunjuk yang diberikan atau hubungi panitia.
	\item \textbf{Reset Akun:} Dalam kasus yang sangat jarang terjadi, seperti salah memberikan akun kepada pemilih atau terjadi kesalahan teknis, panitia memiliki wewenang untuk mereset status pemilihan Anda setelah melalui verifikasi yang ketat.
\end{itemize}

\notebox{Selamat! Anda telah berhasil menggunakan hak suara Anda secara elektronik melalui e-SPPO. Terima kasih atas partisipasi Anda dalam menyukseskan pemilihan yang jujur, adil, dan modern.}

% --- Akhir dari Bab 4: Panduan Penggunaan SSE ---
